% =====================================================================
% Clase -- Álgebra 8° -- Trimestre I -- Semana 04
% Sesiones 010 (mar 22 sep., 100 min), 011 (jue 24 sep., 100 min) y
%          012 (vie 25 sep., 50 min)
% Tema: El escándalo pitagórico (Tema 2 de la guía; sigue en la semana 05)
% Material del docente (no se entrega a los estudiantes).
% =====================================================================
\documentclass[11pt]{article}
\makeatletter\def\input@path{{../../../../../plantillas/estilo/}}\makeatother
\usepackage{mmcantillo}
\guiadelestudiante{../../guia-didactica/guia-periodo-I-numeros-irracionales}

\tipodocumento{Clase}
\asignatura{Álgebra}
\titulo{Semana 4: el escándalo pitagórico}
\grado{8°}
\periodo{I}
% DBA: matematicas grado 8 · DBA 1 — Reconoce la existencia de los números irracionales como
%      números no racionales y los describe de acuerdo con sus características y propiedades.
%      (sesiones 010–012) · DBA 2 — Construye representaciones, argumentos y ejemplos de
%      propiedades de los números racionales y no racionales. (sesión 012)
% Evidencias: DBA 1, representaciones decimales para argumentar por qué un número no es
%      racional (010, 011) y procedimientos geométricos para representar irracionales (012);
%      DBA 2, construir √2 y ubicarlo en la recta, justificando el procedimiento (012).

\begin{document}

\encabezado*

\begin{center}
{\renewcommand{\arraystretch}{1.3}
\begin{tabularx}{\linewidth}{@{}l l l X l@{}}
  \toprule
  \textbf{Sesión} & \textbf{Fecha} & \textbf{Min} & \textbf{Subtema} & \textbf{Entregables}\\
  \midrule
  010 & mar 22 sep., 07:50 & 100 & Hipaso de Metaponto y los inconmensurables & Taller 4 (en clase)\\
  011 & jue 24 sep., 10:10 & 100 & $\sqrt{2}$ no es un número racional & —\\
  012 & vie 25 sep., 11:50 & 50  & Construcción de $\sqrt{2}$: la diagonal del cuadrado & Tarea 4 (para el mar 29 sep.)\\
  \bottomrule
\end{tabularx}}
\end{center}

La semana responde la pregunta que quedó abierta el viernes 18: sí existen decimales que no
son exactos ni periódicos, y el primero que se demostró es la diagonal del cuadrado de lado
$1$. En el hilo del trimestre («La cacería de $\pi$») es la primera presa: demostrar que
$\sqrt{2}$ no es fracción cabe en media página; con $\pi$ costó dos mil años más.
Referencia: \refguia{tema:escandalopitagorico}.

\begin{nota}
Hipaso se cuenta como \emph{leyenda}, no como hecho: no sabemos quién descubrió los
inconmensurables, y ninguna fuente antigua liga directamente a Hipaso con ese descubrimiento
(así lo dice la guía, siguiendo a Huffman). Lo verificable es que en el siglo IV a.\,C. la
demostración ya circulaba: Aristóteles la usa como ejemplo de razonamiento por el absurdo.
\end{nota}

% ---------------------------------------------------------------------
\seccion{Sesión 010 — martes 22 de septiembre · 07:50 · 100 min}

\textbf{Objetivo.} Conocer la convicción pitagórica de que «todo es número» y comprobar, con
fracciones cada vez más cercanas, que ninguna parece dar exactamente la diagonal del cuadrado
de lado $1$.

\textbf{Materiales.} Tablero; regla, escuadra y calculadora; Tarea 3 para revisar; Taller 4
impreso (uno por pareja).

\begin{center}
{\renewcommand{\arraystretch}{1.3}
\begin{tabularx}{\linewidth}{@{}l l X@{}}
  \toprule
  \textbf{Momento} & \textbf{Min} & \textbf{Actividad}\\
  \midrule
  Revisión de la Tarea 3 & 0–15 & Recoger la tarea. Poner en común el reto (punto 4): los
    números que propusieron y la conclusión de que, si nunca repiten un patrón, no pueden ser
    una fracción.\\
  Historia & 15–35 & «Todo es número»: los pitagóricos (siglo VI a.\,C.) creían que toda
    longitud era una razón entre enteros; lo comprobaban con la música (cuerdas en razón
    $2:1$ suenan a una octava). Luego la leyenda del pitagórico que reveló un secreto y murió
    en el mar, al que la tradición llamó Hipaso. Pregunta al grupo: «¿qué secreto podía ser
    tan grave?»\\
  Exploración & 35–55 & Cada pareja dibuja un cuadrado de $10$ cm de lado, mide la diagonal
    (unos $14{,}1$ cm) y busca una fracción del lado que la dé. Recoger propuestas en el
    tablero ($\frac{7}{5}$, $\frac{141}{100}$…) y elevar al cuadrado: ninguna da $2$. Se
    anuncia que el jueves se demostrará que \emph{ninguna} puede darlo.\\
  Taller 4 & 55–95 & En parejas. Circula por los puestos; el punto 1 es el central (las
    fracciones se acercan a $\sqrt{2}$ sin alcanzarla).\\
  Cierre & 95–100 & Recoger el taller. Pregunta para el jueves: «¿cómo se puede estar seguro
    de que \emph{ninguna} fracción funciona, si hay infinitas fracciones?»\\
  \bottomrule
\end{tabularx}}
\end{center}

\begin{nota}
Los puntos 3 y 4 del taller ubican irracionales en la recta por su aproximación decimal;
basta con que los ubiquen entre dos enteros consecutivos. La ubicación exacta con compás es
de la sesión 012 y de la semana 05. Si el tiempo no alcanza, el punto 4 queda para la casa.
\end{nota}

% ---------------------------------------------------------------------
\seccion{Sesión 011 — jueves 24 de septiembre · 10:10 · 100 min}

\textbf{Objetivo.} Seguir y reconstruir la demostración por el absurdo de que $\sqrt{2}$ no es
racional, y definir número irracional.

\begin{center}
{\renewcommand{\arraystretch}{1.3}
\begin{tabularx}{\linewidth}{@{}l l X@{}}
  \toprule
  \textbf{Momento} & \textbf{Min} & \textbf{Actividad}\\
  \midrule
  Enganche & 0–10 & Retomar la pregunta del martes: probar fracción por fracción no
    termina nunca; hace falta otra forma de razonar.\\
  Preparación & 10–30 & Pares e impares: todo par es $2k$; el cuadrado de un impar es impar
    ($3^2, 5^2, 7^2$…) y el de un par es par. Conclusión que se usará: si $p^2$ es par,
    entonces $p$ es par. Fracción reducida: sin factores comunes.\\
  Demostración & 30–60 & En el tablero, paso a paso (la de la guía, Teorema «$\sqrt{2}$ es
    irracional»): suponer $\sqrt{2} = \frac{p}{q}$ reducida; $p^2 = 2q^2$; $p$ par, $p = 2k$;
    $q^2 = 2k^2$; $q$ par; contradicción. Nombrar el método: \emph{por el absurdo}. Leer la
    frase de Aristóteles: los impares resultarían iguales a los pares.\\
  Reconstrucción & 60–85 & En parejas, sin mirar el tablero, escriben la demostración en el
    cuaderno con sus palabras; luego la comparan con la de la guía y marcan el paso que les
    faltó.\\
  Definición y cierre & 85–100 & Definición de número irracional (decimal infinito no
    periódico; conjunto $\mathbb{I}$) y las tres caras de $\sqrt{2}$: la diagonal, el decimal
    $1{,}41421356\ldots$ y el número positivo cuyo cuadrado es $2$. Pregunta para discutir:
    «¿dónde falla la demostración si la intentamos con $\sqrt{4}$?» ($4 = 2^2$: no hay
    contradicción, y en efecto $\sqrt{4} = 2$).\\
  \bottomrule
\end{tabularx}}
\end{center}

\begin{nota}
El paso que más cuesta es «si $p^2$ es par, $p$ es par»: justifícalo con la tabla de la
preparación (un impar al cuadrado siempre da impar), no lo des por evidente. La práctica natural
de esta sesión sería adaptar la demostración a $\sqrt{3}$ (ejercicio \texttt{irracionales-8-009}
del banco), pero su respuesta modelo sigue pendiente de aprobación: úsalo solo si ya lo
aprobaste.
\end{nota}

% ---------------------------------------------------------------------
\seccion{Sesión 012 — viernes 25 de septiembre · 11:50 · 50 min}

\textbf{Objetivo.} Construir $\sqrt{2}$ con regla y compás como la diagonal del cuadrado de
lado $1$, llevarla a la recta numérica y justificar por qué el punto obtenido es exactamente
$\sqrt{2}$.

\textbf{Materiales.} Regla, compás y escuadra (pedirlos el jueves).

\begin{center}
{\renewcommand{\arraystretch}{1.3}
\begin{tabularx}{\linewidth}{@{}l l X@{}}
  \toprule
  \textbf{Momento} & \textbf{Min} & \textbf{Actividad}\\
  \midrule
  Construcción guiada & 0–20 & Con unidad de $3$ cm: recta numérica, cuadrado de lado $1$
    sobre el segmento $[0,1]$, diagonal desde $0$; compás en $0$ con abertura hasta el
    extremo de la diagonal; arco hasta la recta. Marcar $\sqrt{2}$ y comprobar con la regla
    que queda cerca de $1{,}41$ unidades. La diagonal mide $\sqrt{1^2 + 1^2}$ por el teorema de
    Pitágoras, que se usa como herramienta (se estudia en Geometría).\\
  Justificación & 20–30 & «¿Por qué el punto es \emph{exactamente} $\sqrt{2}$ y no
    $1{,}41$?» Porque el compás conserva la longitud: el arco lleva la diagonal entera, no una
    medida con la regla.\\
  Segundo caso & 30–42 & Con el triángulo de catetos $1$ y $2$, llevar $\sqrt{5}$ a la recta
    (ejemplo de la guía, \texttt{irracionales-8-010}); también $-\sqrt{2}$, con el arco hacia la
    izquierda.\\
  Cierre y tarea & 42–50 & Anunciar la espiral de Teodoro (semana 05) y asignar la Tarea 4,
    para el martes 29 de septiembre.\\
  \bottomrule
\end{tabularx}}
\end{center}

\begin{nota}
Teorema de Pitágoras solo con el cuadrado de lado $1$ y el triángulo de catetos $1$ y $2$
(decisión del plan de la guía, hasta la semana 09). Si sobra tiempo, los estudiantes pueden
empezar la pregunta de la espiral de Teodoro de la guía, que se trabaja en la semana 05.
\end{nota}

\end{document}
